\section{Coding Performance and Image Quality}
\label{sec:metrics}

Images compressed using lossless compression can be perfectly reconstructed by the decompressor. Near-lossless compression introduces noise in the reconstructed image, due imperfect reconstruction. The error limits of the issue 2 in-loop quantizer can be tuned, adjusting the maximum error $m_z(t)$ of reconstructed samples. As described by \citeauthor{salomonDataCompressionComplete2004} in, compression noise can be characterized using \acrfull{PSNR}, giving an indication as to the quality to the reconstructed image compared to the original \cite{salomonDataCompressionComplete2004}. For a single spectral band $z$, the \acrshort{PSNR} is given by

\begin{equation}
	\mathrm{PSNR}_z = 10 \cdot \text{log}_{10} \left( \frac{\mathrm{MAX}^2_I}{\mathrm{MSE_z}} \right),
	\label{eq:PSNRz}
\end{equation}

where the \acrfull{MSE} of the spectral band is given by

\begin{equation}
	\mathrm{MSE_z} = \frac{1}{N_x N_y}
	\sum_{i=0}^{N_x-1} \sum_{j=0}^{N_y-1} \bigl[I_z(i,j) - K_z(i,j)\bigr]^2.
	\label{eq:MSE}
\end{equation}

Samples corresponding to the spectral bands of the original and decompressed image are denoted using $I_z$ and $K_z$. The maximum sample value of the original image is denoted $\mathrm{MAX}_I$. The \acrshort{PSNR} for the entire image is found by averaging over $\mathrm{PSNR}_z$ as

\begin{equation}
	\mathrm{PSNR} = \frac{1}{N_z} \sum_{z=0}^{N_z-1} \mathrm{PSNR}_z.
	\label{eq:PSNR}
\end{equation}

Evaluating the compression performance is done in terms of \acrfull{CR} and Bits Per Sample (bitrate), given by

\begin{equation}
	\mathrm{CR} = \frac{N_x \cdot N_y \cdot N_z \cdot D}{\text{compressed data size (bits)}},
	\label{eq:cr}
\end{equation}

and

\begin{equation}
	\text{Bits Per Sample} = \frac{D}{\mathrm{CR}},
	\label{eq:bitrate}
\end{equation}

where dynamic range $D$ is the number of bits per sample.

